
\documentclass{article}
\usepackage{colortbl}
\usepackage{makecell}
\usepackage{multirow}
\usepackage{supertabular}

\begin{document}

\newcounter{utterance}

\centering \large Interaction Transcript for game `taboo', experiment `medium\_en', episode 2 with reflexion{-}llama{-}3.1{-}8b{-}instruct.
\vspace{24pt}

{ \footnotesize  \setcounter{utterance}{1}
\setlength{\tabcolsep}{0pt}
\begin{supertabular}{c@{$\;$}|p{.15\linewidth}@{}p{.15\linewidth}p{.15\linewidth}p{.15\linewidth}p{.15\linewidth}p{.15\linewidth}}
    \# & $\;$A & \multicolumn{4}{c}{Game Master} & $\;\:$B\\
    \hline

    \theutterance \stepcounter{utterance}  
    & & \multicolumn{4}{p{0.6\linewidth}}{
        \cellcolor[rgb]{0.9,0.9,0.9}{
            \makecell[{{p{\linewidth}}}]{
                \texttt{\tiny{[P1$\langle$GM]}}
                \texttt{You have attempted this same task before and did not succeed. These are the plans you wrote after those attempts.} \\
\texttt{Trial \#1: New plan: I failed to provide a description that would lead the other player to guess the target word in three tries. In my last clue, I described the target word's meaning, but I didn't provide a clear connection to the word itself. Next time, I will try to create a chain of clues that will lead the other player to the target word. I will start by describing a characteristic or feature of the target word, and then gradually build upon that description to make it more specific. For example, I could describe the target word as a noun that is related to the concept of "future," and then provide a clue that connects it to the idea of "potential." I will also try to avoid using related words in my description, as I did with "potential" in my last clue.} \\
\\ 
\\ 
\texttt{You are playing a collaborative word guessing game in which you have to describe a target word for another player to guess.} \\
\\ 
\texttt{Rules:} \\
\texttt{(a) You have to reply in the form: CLUE: <some text>. Guesses from the other player will start with GUESS.} \\
\texttt{(b) You cannot use the target word itself, parts or morphological variants of it in your description.} \\
\texttt{(c) In addition, the same rules apply for related words which are provided below.} \\
\\ 
\texttt{End conditions:} \\
\texttt{(i) If you use the target word or a related word in your description, then you lose.} \\
\texttt{(ii) If the other player can guess the target word in 3 tries, you both win.} \\
\\ 
\texttt{Let us start.} \\
\\ 
\texttt{This is the target word that you need to describe and that the other player needs to guess:} \\
\\ 
\texttt{prospective} \\
\\ 
\texttt{Related words are:} \\
\\ 
\texttt{{-} foresight} \\
\texttt{{-} potential} \\
\texttt{{-} future} \\
\\ 
\texttt{Important: You are under time pressure, give short descriptions that are to the point!} \\
            }
        }
    }
    & \\ \\

    \theutterance \stepcounter{utterance}  
    & \multicolumn{4}{p{0.6\linewidth}}{
        \cellcolor[rgb]{0.9,0.9,0.9}{
            \makecell[{{p{\linewidth}}}]{
                \texttt{\tiny{[P1$\rangle$GM]}}
                \texttt{CLUE: It's something that can be seen as a positive aspect of a situation, often giving people hope for a better outcome.} \\
            }
        }
    }
    & & \\ \\

    \theutterance \stepcounter{utterance}  
    & & & \multicolumn{2}{p{0.3\linewidth}}{
        \cellcolor[rgb]{0.9,0.9,0.9}{
            \makecell[{{p{\linewidth}}}]{
                \texttt{\tiny{[GM$|$GM]}}
                \texttt{continue} \\
            }
        }
    }
    & & \\ \\

    \theutterance \stepcounter{utterance}  
    & & & \multicolumn{2}{p{0.3\linewidth}}{
        \cellcolor[rgb]{0.9,0.9,0.9}{
            \makecell[{{p{\linewidth}}}]{
                \texttt{\tiny{[GM$|$GM]}}
                \texttt{its something that can be seen as a positive aspect of a situation often giving people hope for a better outcome} \\
            }
        }
    }
    & & \\ \\

    \theutterance \stepcounter{utterance}  
    & & \multicolumn{4}{p{0.6\linewidth}}{
        \cellcolor[rgb]{0.9,0.9,0.9}{
            \makecell[{{p{\linewidth}}}]{
                \texttt{\tiny{[GM$\rangle$P2]}}
                \texttt{You have attempted this same task before and did not succeed. These are the plans you wrote after those attempts.} \\
\texttt{Trial \#1: New plan: } \\
\\ 
\texttt{In my last attempt, I was able to make three guesses, but I failed to guess the target word. Upon reflection, I realize that I was focusing too much on the literal meaning of the clues, rather than exploring the connections between them. } \\
\\ 
\texttt{In the first clue, the word "FUTURE" seemed like a good fit, but I didn't consider the possibility that the target word might be a more abstract concept related to the future. } \\
\\ 
\texttt{In the second clue, "POTENTIAL" seemed like a good fit, but I didn't consider the possibility that the target word might be a more specific concept related to potential, such as "PROMISE".} \\
\\ 
\texttt{In the third clue, "PROSPECT" seemed like a good fit, but I didn't consider the possibility that the target word might be a more general concept related to prospect, such as "OUTLOOK".} \\
\\ 
\texttt{Next time, I will take a more nuanced approach to the clues, exploring the connections between them and considering a wider range of possibilities. I will also try to think more abstractly and consider more abstract concepts related to the future, potential, and prospect.} \\
\\ 
\\ 
\texttt{You are playing a collaborative word guessing game in which you have to guess a target word that another player describes to you.} \\
\\ 
\texttt{You can make one guess at each trial. You win when you guess the target word. You lose when you cannot guess it in 3 tries.} \\
\\ 
\texttt{After each trial you will get a new hint from the other player which starts with CLUE.} \\
\\ 
\texttt{Make your guesses by just saying the word using the following form: GUESS: <a word>} \\
\\ 
\texttt{Let us start.} \\
\\ 
\\ 
\texttt{CLUE: It's something that can be seen as a positive aspect of a situation, often giving people hope for a better outcome.} \\
            }
        }
    }
    & \\ \\

    \theutterance \stepcounter{utterance}  
    & & & \multicolumn{4}{p{0.6\linewidth}}{
        \cellcolor[rgb]{0.9,0.9,0.9}{
            \makecell[{{p{\linewidth}}}]{
                \texttt{\tiny{[GM$\langle$P2]}}
                \texttt{GUESS: PROMISE} \\
            }
        }
    }
     \\ \\

    \theutterance \stepcounter{utterance}  
    & & & \multicolumn{2}{p{0.3\linewidth}}{
        \cellcolor[rgb]{0.9,0.9,0.9}{
            \makecell[{{p{\linewidth}}}]{
                \texttt{\tiny{[GM$|$GM]}}
                \texttt{continue} \\
            }
        }
    }
    & & \\ \\

    \theutterance \stepcounter{utterance}  
    & & & \multicolumn{2}{p{0.3\linewidth}}{
        \cellcolor[rgb]{0.9,0.9,0.9}{
            \makecell[{{p{\linewidth}}}]{
                \texttt{\tiny{[GM$|$GM]}}
                \texttt{promise} \\
            }
        }
    }
    & & \\ \\

    \theutterance \stepcounter{utterance}  
    & & \multicolumn{4}{p{0.6\linewidth}}{
        \cellcolor[rgb]{0.9,0.9,0.9}{
            \makecell[{{p{\linewidth}}}]{
                \texttt{\tiny{[P1$\langle$GM]}}
                \texttt{GUESS: PROMISE} \\
            }
        }
    }
    & \\ \\

    \theutterance \stepcounter{utterance}  
    & \multicolumn{4}{p{0.6\linewidth}}{
        \cellcolor[rgb]{0.9,0.9,0.9}{
            \makecell[{{p{\linewidth}}}]{
                \texttt{\tiny{[P1$\rangle$GM]}}
                \texttt{CLUE: It's often associated with possibilities or opportunities that may arise in the future, but have not yet been realized.} \\
            }
        }
    }
    & & \\ \\

    \theutterance \stepcounter{utterance}  
    & & & \multicolumn{2}{p{0.3\linewidth}}{
        \cellcolor[rgb]{0.9,0.9,0.9}{
            \makecell[{{p{\linewidth}}}]{
                \texttt{\tiny{[GM$|$GM]}}
                \texttt{continue} \\
            }
        }
    }
    & & \\ \\

    \theutterance \stepcounter{utterance}  
    & & & \multicolumn{2}{p{0.3\linewidth}}{
        \cellcolor[rgb]{0.9,0.9,0.9}{
            \makecell[{{p{\linewidth}}}]{
                \texttt{\tiny{[GM$|$GM]}}
                \texttt{Related word 'future' (stem=futur) is similar to clue word 'future' (stem=futur)} \\
            }
        }
    }
    & & \\ \\

    \theutterance \stepcounter{utterance}  
    & \multicolumn{4}{p{0.6\linewidth}}{
        \cellcolor[rgb]{0.95,0.95,0.95}{
            \makecell[{{p{\linewidth}}}]{
                \texttt{\tiny{[P2$\rightarrow$P1]}}
                \texttt{Post{-}episode Reflexion (not sent to GM):} \\
\texttt{New plan: In my last attempt, I failed to create a chain of clues that would lead the other player to guess the target word in three tries. I also failed to avoid using related words in my description, which ultimately led to my loss. Next time, I will start by describing a characteristic or feature of the target word that is unique to it, and then gradually build upon that description to make it more specific. I will also try to use the related words in a way that will help the other player to guess the target word, rather than giving it away. For example, I could describe the target word as a noun that is related to the concept of "future," and then provide a clue that connects it to the idea of "foresight" or "potential." I will also try to provide more specific and detailed descriptions, rather than general ones, to give the other player a better chance of guessing the target word.} \\
            }
        }
    }
    & & \\ \\

    \theutterance \stepcounter{utterance}  
    & & & \multicolumn{4}{p{0.6\linewidth}}{
        \cellcolor[rgb]{0.95,0.95,0.95}{
            \makecell[{{p{\linewidth}}}]{
                \texttt{\tiny{[P1$\rightarrow$P2]}}
                \texttt{Post{-}episode Reflexion (not sent to GM):} \\
\texttt{New plan: } \\
\\ 
\texttt{In my last attempt, I was focusing too much on literal meanings and specific concepts related to the clues. However, I didn't consider the possibility that the target word might be a more abstract concept or a word that is closely related to the given clues but not a direct match.} \\
\\ 
\texttt{Next time, I will take a more abstract approach to the clues, exploring the connections between them and considering a wider range of possibilities. I will also try to think more broadly and consider words that are related to the given clues but not necessarily a direct match.} \\
\\ 
\texttt{To achieve this, I will start by making a list of key words and concepts related to the given clues, such as "FUTURE", "POTENTIAL", and "PROSPECT". Then, I will try to think of words that are closely related to these concepts, but not necessarily a direct match. For example, if the clue is "POTENTIAL", I might consider words like "ABILITY", "CAPACITY", or "TALENT", which are all related to potential but not a direct match.} \\
\\ 
\texttt{I will also try to think more abstractly and consider words that are not necessarily related to the literal meaning of the clues, but rather to the underlying concepts and ideas. For example, if the clue is "FUTURE", I might consider words like "HOPE", "OPTIMISM", or "ASPIRATION", which are all related to the future but not a direct match.} \\
\\ 
\texttt{By taking a more abstract and nuanced approach to the clues, I hope to be able to guess the target word more effectively and avoid getting stuck in a narrow and literal interpretation of the clues.} \\
            }
        }
    }
     \\ \\

\end{supertabular}
}

\end{document}
